\section{Introduction}

Code obfuscation is widely used to make software difficult to inspect
while preserving its intended functionality. In benign settings,
developers employ obfuscation to protect intellectual property,
implementation details, and proprietary business logic. In adversarial
settings, however, malware authors routinely exploit the same
techniques to conceal malicious behavior and frustrate program
analysis and reverse engineering. Common transformations include
identifier renaming, control-flow restructuring, dead-code insertion,
expression rewriting, call indirection, etc. Although such
transformations are intended to preserve program semantics, they can
substantially obscure how those semantics are expressed in the source
code, thereby increasing the difficulty of understanding the
behavior of malware code~\cite{collberg1997taxonomy,collberg2002manufacturing,
Obfuscation-The-Hidden-Malware,
An-Observational-Investigation-of-Reverse-Engineers-Process-and-Mental-Models}.

This difficulty is well established for human analysts. Nguyen
{\em et al.}~\cite{nguyen2026effectcodeobfuscationhuman}, for example,
showed that code obfuscation significantly reduces humans' accuracy in
reasoning about program outputs. Such findings are particularly
important for malware analysis, where analysts must often infer the
behavior of code whose surface structure has been deliberately
engineered to resist comprehension.

Large language models (LLMs) offer a promising alternative for
automating such reasoning. Models have demonstrated strong
capabilities in code generation, summarization, vulnerability
detection, and program understanding
\cite{chen2021evaluatinglargelanguagemodels,li2022codebert,
guo2021graphcodebert,roziere2023codegen,
Competition-level-code-generation-with-AlphaCode}.
However, most evaluations are performed on naturally written programs,
where meaningful identifiers, familiar control structures, and common
coding idioms provide abundant lexical and structural cues. Obfuscated
programs remove or distort many of these cues while retaining the
underlying computation, creating a useful stress test of whether an
LLM actually reasons about program behavior or instead depends on
surface regularities.

Recent studies suggest that this distinction matters.
Rong {\em et al.}~\cite{ase26} showed that semantics-preserving
obfuscations can cause substantial degradation in LLM program
reasoning, including program-output prediction. Similarly, Nikiema
{\em et al.}~\cite{nikiema2025codebarrierllmsactually} reported a
statistically significant decline in model performance as the
complexity of obfuscation increases. Together, these findings indicate
that this limitation is especially problematic
for security applications: an AI-assisted malware analyzer should not
cease to understand a malicious computation merely due to obfuscation.

%today's code LLMs are not invariant to transformations that leave program semantics unchanged.

An important challenge, however, has received considerably less
attention: \emph{real obfuscated programs rarely contain only one
transformation}. Obfuscators can combine multiple techniques so that,
for example, renamed identifiers are embedded inside flattened control
flow, rewritten expressions, and additional irrelevant code. Such
\emph{stacked obfuscation} creates a fundamentally harder deployment
setting. Success on each transformation independently does not
necessarily imply success when those transformations occur together.
Consequently, simply demonstrating that a model can be adapted to
individual obfuscations is insufficient. For practical malware
analysis, the learned capability must also \emph{compose}.

This observation naturally suggests several ways to compose knowledge
learned from individual transformations. One could train a specialist
for each obfuscation and use a \emph{router} to select the appropriate
specialist at inference time. Alternatively, one could combine the
specialists in parameter space through \emph{model merging}. A third,
and seemingly more direct, strategy is to increase \emph{training
breadth} by fine-tuning a single model on examples from many individual
obfuscations, hoping that knowledge learned separately will recombine
when transformations are stacked. We systematically evaluate these
three strategies and find that none provides the desired robustness.
The learned router performs no better than a random routing policy;
merging specialized models performs at or below the clean-code-tuned
control; and broader training, although effective on stacks assembled
from transformations represented during training, introduces a
substantial performance penalty on an unseen obfuscation
family. Thus, broader exposure can improve recombination within the
training distribution while simultaneously making the model
\emph{less} robust outside that distribution.

Our analysis reveals why. The apparent robustness obtained from
obfuscation-specific adaptation does not primarily arise from learning
the semantics that remain invariant under transformation. Instead,
models learn regularities associated with the \emph{surface artifacts}
left by particular obfuscators. Evidence for this behavior appears
across multiple analyses: adaptation transfers poorly between distinct
mechanisms within the same obfuscation family; the benefit of broad
training is substantially stronger when recognizable identifier cues
remain available; and improvements disappear when the same programs
are queried in a different reasoning direction. Thus, stacking exposes
a deeper problem than merely insufficient training coverage:
obfuscation-specific fine-tuning can teach the model how particular
transformations \emph{look} without teaching it to consistently recover
the behavior that those transformations preserve.

This diagnosis suggests a different principle for robust adaptation.
Rather than attempting to compose transformation-specific knowledge,
we anchor learning to something that is invariant across all such
transformations: the model's behavior on the corresponding
\emph{unobfuscated program}. For every obfuscated training program
$x_{\mathrm{obf}}$, we retain its clean parent
$x_{\mathrm{clean}}$. A teacher model, tuned for reasoning over clean
code, processes $x_{\mathrm{clean}}$, while the student processes
$x_{\mathrm{obf}}$. In addition to the ordinary task loss, we minimize
the KL divergence between their output distributions:
%
\begin{equation}
\mathcal{L}
=
\mathrm{CE}\!\left(y \mid x_{\mathrm{obf}}\right)
+
\lambda\,
D_{\mathrm{KL}}\!\left(
p_{\mathrm{teacher}}(\cdot \mid x_{\mathrm{clean}})
\,\Vert\,
p_{\mathrm{student}}(\cdot \mid x_{\mathrm{obf}})
\right).
\label{eq:paired_consistency}
\end{equation}
%
The first term exposes the student to obfuscated code, whereas the
second explicitly encourages it to preserve the behavioral response of
a model reasoning over the semantics-equivalent clean program.
Importantly, this objective does not require recognizing which
obfuscations are present, selecting among specialists, merging
transformation-specific parameters, or recovering the original source
code. Instead, it
changes the learning target itself: rather than anchoring adaptation to
the surface patterns of individual transformations, it anchors the
student to the clean-code behavior that those transformations should
preserve.

Our empirical results demonstrate the advantage of this
\emph{paired-consistency anchoring}. In contrast to routing, model merging, and breadth-only training, the KL-based objective retains the
benefits of learning from diverse obfuscations without paying the
corresponding penalty on an unseen obfuscation family. At the 7B scale,
the anchored model improves over breadth training by 4.59 percentage
points on the unseen family while introducing no degradation on clean
code. Additional stress tests show that the resulting model continues
to perform well under deeper compositions and transfers to a second
programming language. At the same time, our analysis carefully bounds
the claim: the method does not create a new ability to solve arbitrary
unseen obfuscations beyond that of the clean-code teacher. Rather, it
prevents obfuscation adaptation from destroying capabilities the model
already possesses. This distinction is important: robust reasoning
over obfuscated programs requires not only acquiring adaptation gains,
but also preserving semantic competence as the surface form of the
program changes.

{\em Contributions.} This paper makes the following main contributions:

\textbf{1. A systematic study and empirical evaluation of compositional robustness to
stacked obfuscation.}
    We investigate whether capabilities learned from individual
    semantics-preserving transformations compose when multiple
    transformations are stacked, a setting representative of practical
    obfuscation and malware analysis.

    %\item \textbf{An empirical evaluation of three natural composition strategies.}
    %We study inference-time routing, weight-space model merging, and broad multi-transformation training. We show that routing provides essentially no advantage over random selection, merging fails to combine specialist capabilities, and breadth training improves reasoning over seen transformation combinations but significantly degrades performance on an unseen obfuscation family.

\textbf{2. An empirical characterization of why adaptation
    fails to compose.}
    With complementary analyses across transformation mechanisms,
    identifier dependence, reasoning direction, and confidence
    behavior, we show that obfuscation adaptation predominantly learns
    transformation-specific surface regularities rather than a general
    invariance to semantics-preserving transformations.

\textbf{3. A clean-behavior anchoring objective for robust
    obfuscation adaptation.}
    We introduce a paired-consistency training objective that minimizes
    the KL divergence between a clean-code teacher and an
    obfuscated-code student. The approach preserves the gains of broad
    obfuscation training while eliminating its degradation on unseen
    transformation families and clean code, outperforming routing,
    merging, and breadth-only adaptation.

\textbf{4. Stress testing and characterization of the resulting
    robustness.} We evaluate the anchored model under deeper obfuscation stacks, a
    second programming language, unseen transformations, and reversed
    reasoning tasks. 
    
    %These experiments clarify both the benefits and the limits of the proposed approach: anchoring preserves the teacher's semantic reasoning capability under obfuscation rather than creating unsupported new invariance.

%\end{itemize}